\section{Group configurations}
\label{sec:gc}

\subsection{Definition and first consequences}

\begin{definition}[group configuration]\label{def:gc}
A \emph{group configuration over $k$} is a $6$-tuple $(a,b,c,x,y,z)$ of tuples of
$K$ such that, with the \emph{lines}
\[
   L_1=\{a,b,c\},\quad L_2=\{a,x,y\},\quad L_3=\{b,y,z\},\quad L_4=\{c,z,x\},
\]
the following hold:
\begin{enumerate}[label=\textup{(GC\arabic*)}]
\item\label{GC:group} $\acl(ab)=\acl(ac)=\acl(bc)$;
\item\label{GC:space} $\acl(ax)=\acl(ay)$, $\acl(by)=\acl(bz)$, $\acl(cx)=\acl(cz)$;
\item\label{GC:indep} every triple of points not contained in a line is independent:
  $\rk(pqr)=\rk(p)+\rk(q)+\rk(r)$ for all non-collinear $\{p,q,r\}$;
\item\label{GC:cb} $a\subseteq\acl(\Cb(xy/a))$, $b\subseteq\acl(\Cb(yz/b))$, $c\subseteq\acl(\Cb(xz/c))$;
\item\label{GC:rank} $\rk(a)\ge1$.
\end{enumerate}
The points $a,b,c$ are the \emph{group points}, $x,y,z$ the \emph{space points}.
\end{definition}

Two distinct lines meet in exactly one point ($L_1\cap L_2=\{a\}$, $L_1\cap L_3=\{b\}$,
$L_1\cap L_4=\{c\}$, $L_2\cap L_3=\{y\}$, $L_2\cap L_4=\{x\}$, $L_3\cap L_4=\{z\}$), each
point lies on exactly two lines, and any two points lie on at most one line (the three
pairs $\{a,z\}$, $\{b,x\}$, $\{c,y\}$ lie on no line).  The group line $L_1$ is
distinguished by the definition (\ref{GC:group} and \ref{GC:cb} concern group points
only), and the definition is invariant under the six relabelings induced by
permutations of the three space lines $L_2,L_3,L_4$: these permute $a,b,c$ and
correspondingly $x,y,z$ (for instance $L_3\leftrightarrow L_4$ is $b\leftrightarrow c$,
$y\leftrightarrow x$, $z$ fixed), and each of \ref{GC:group}--\ref{GC:cb} is symmetric in
the two space points of a line and in the three group points (\ref{GC:cb} because $\FOD$
is invariant under renaming variables, so $\Cb(xy/a)=\Cb(yx/a)$); \ref{GC:rank} is
invariant because $\rk(a)=\rk(b)=\rk(c)$ by \cref{lem:gc-ranks}.

\begin{lemma}[ranks]\label{lem:gc-ranks}
Let $(a,b,c,x,y,z)$ be a group configuration.  Then $\rk(a)=\rk(b)=\rk(c)=:n$,
$\rk(x)=\rk(y)=\rk(z)=:m$, $1\le m\le n$; all pairs of points are independent;
$\rk(abc)=2n$; $\rk(axy)=\rk(byz)=\rk(cxz)=n+m$; $\rk(abcxyz)=2n+m$; and each point
is independent from each line not through it: for instance $a\ind{}byz$, $a\ind{}cxz$,
$b\ind{}axy$, $b\ind{}cxz$, $c\ind{}axy$, $c\ind{}byz$, $x\ind{}abc$, $x\ind{}byz$, and
so on.  Moreover $m\ge 1$ follows from $n\ge1$.
\end{lemma}
\begin{proof}
Every pair lies in a non-collinear triple (a triple containing a given pair and lying
on a line lies on the unique line through that pair, if any; so any third point off that
line works), so pairs are independent by \ref{GC:indep}.  By \ref{GC:group}, $\rk(ab)=\rk(ac)$, i.e.\ $n_b=n_c$,
and similarly all three agree.  By \ref{GC:space}, $\rk(ax)=\rk(ay)$ so $\rk(x)=\rk(y)$,
etc.  $\rk(axy)=\rk(ay)=n+m\ge\rk(xy)=2m$ gives $m\le n$.  For $\rk(abcxyz)$: $(a,b,y)$ is
non-collinear, so $\rk(aby)=2n+m$, and $c\in\acl(ab)$, $x\in\acl(ay)$, $z\in\acl(by)$.
Independence of a point $p$ from a line $\ell$ not through $p$: pick two points $q,r$
of $\ell$, taking $q$ to be the group point of $\ell$ if $\ell$ is a space line; then
$(p,q,r)$ is non-collinear (a line containing $q$ and $r$ is $\ell$, as two distinct
lines share only one point, and $p\notin\ell$), so $p\ind{}qr$, and the third point of $\ell$ is algebraic over $qr$ by \ref{GC:group} (if
$\ell=L_1$) or \ref{GC:space} (a space point is algebraic over the group point and the
other space point of its line); hence $p\ind{}\ell$.  E.g.\ $a\ind{}by$ and
$z\in\acl(by)$ give $a\ind{}byz$.  If $m=0$
then $x,y\subseteq k$, $P_{\acl(a)}(xy)$ is a $k$-rational point and $\Cb(xy/a)=k$,
contradicting \ref{GC:cb} and $n\ge1$.
\end{proof}

\begin{lemma}[invariance]\label{lem:gc-invariance}
Let $\mathfrak c=(a,b,c,x,y,z)$ be a group configuration over $k$.
\begin{enumerate}[nosep]
\item For $\sigma\in\Gamma$, $\sigma\mathfrak c$ is a group configuration over $k$.
\item (base change) If $e$ is a tuple with $e\ind{}abcxyz$, then $\mathfrak c$ is a group
  configuration over $k'=\acl(k,e)$ (all notions of
  \cref{def:closures,def:cb} now taken relative to the base $k'$; $K$ still has
  infinite transcendence degree over $k'$).
\item (interalgebraic replacement) If $a'\interalg a,\dots,z'\interalg z$ then
  $(a',b',c',x',y',z')$ is a group configuration over $k$.
\end{enumerate}
\end{lemma}
\begin{proof}
(1) \cref{lem:equiv} and \ref{Cb:equiv}.
(2) \ref{GC:group}, \ref{GC:space}: e.g.\ $c\in\acl(ab)\subseteq\acl(k'ab)$ and
symmetrically.  \ref{GC:indep}: for a non-collinear triple, $\rk(pqr/e)=\rk(pqr)$ and
$\rk(p/e)=\rk(p)$ etc.\ by \ref{R:basechange}.  \ref{GC:cb}: the canonical base over the
new base is $\Cb'(xy/a)=k'\cdot\FOD(P_{\acl(k'a)}(xy))$ and
$\FOD(P_{\acl(ea)}(xy))=\FOD(P_{\acl(a)}(xy))$ by \ref{Cb:basechange} (as
$e\ind{}xya$); hence $\acl(\Cb'(xy/a))=\acl(k',\Cb(xy/a))\supseteq\acl(\Cb(xy/a))\ni a$.
\ref{GC:rank}: $\rk(a/e)=\rk(a)$.
(3) \ref{GC:group}--\ref{GC:indep} and \ref{GC:rank} are invariant under replacing
points by interalgebraic ones (\ref{R:aclinv}).  For \ref{GC:cb}:
$\Cb(x'y'/a')=\Cb(x'y'/\acl(a))$ and $\acl(\Cb(x'y'/\acl a))=\acl(\Cb(xy/\acl a))$ by
\ref{Cb:interalg}, which contains $a$, hence $a'$.
\end{proof}

\begin{lemma}[conjugate sets of a space point]\label{lem:conjset-gc}
Let $(a,b,c,x,y,z)$ be a group configuration and let $r$ be a space point with the two
lines $\{p_1,q_1,r\}$ and $\{p_2,q_2,r\}$ through it, $p_1,p_2$ being their group points
and $q_1,q_2$ their other space points.  Put $L=k(p_1q_1p_2q_2)$.  Then
$r\subseteq\acl(L)$, every element of the conjugate set $\Gamma_L\cdot r$ lies in
$\acl(r)$, and the code $\tilde r=\code{\Gamma_L r}$ satisfies $\tilde r\subseteq\dcl(L)$
and $\tilde r\interalg r$.
\end{lemma}
\begin{proof}
$r\in\acl(p_1q_1)\cap\acl(p_2q_2)$ by \ref{GC:space} (a space point is algebraic over
the group point and the other space point of each line through it).  The group points
$p_1,p_2$ lie on different lines through $r$, and the line through $p_1,p_2$ is $L_1$,
which does not contain $r$; so $(r,p_1,p_2)$ is non-collinear, $p_1\ind{}rp_2$ and hence
$p_1\ind{r}p_2$.  Since $q_1\in\acl(rp_1)$ and $q_2\in\acl(rp_2)$ (\ref{GC:space} again),
\ref{R:aclinv} gives $p_1q_1\ind{r}p_2q_2$, and \ref{R:meet} gives
$\acl(rp_1q_1)\cap\acl(rp_2q_2)=\acl(r)$.  For $\sigma\in\Gamma_L$ we have
$\sigma r\in\sigma\acl(p_1q_1)\cap\sigma\acl(p_2q_2)=\acl(p_1q_1)\cap\acl(p_2q_2)\subseteq\acl(r)$.
The claims about $\tilde r$ are \cref{lem:conjcode}.  (Note that a group point need
\emph{not} be algebraic over the two space points of its line; only the space-point
clauses of \ref{GC:space} were used.)
\end{proof}

\begin{remark}
For a group point the same statement is false in general (cf.\ \cref{rem:cb-not-acl});
this is why the definability step of \cref{sec:proof} only ever codes space points.
\end{remark}

\subsection{Configurations of a homogeneous space}

\begin{proposition}[the converse]\label{prop:converse}
Let $\mathcal G$ be a generically transitive generic set of germs on $V$ closed under
inverse and generic composition, with group $G$ (\cref{thm:hw}), $n=\rk(\mathcal G)$,
$m=\dim V\ge1$.  Let $g,h\in\mathcal G$ and $y\in V_{\Gen}$ with $(g,h,y)$ independent,
and put $x=g(y)$, $z=h(y)$, $c=hg^{-1}\in G$.  Then
\[
  \bigl(\code g,\ \code h,\ \code c,\ x,\ y,\ z\bigr)
\]
is a group configuration over $k$ with $\rk(\code g)=n$ and $\rk(y)=m$; moreover
$z=c(x)$ and $h=cg$.
\end{proposition}
\begin{proof}
Write $a=\code g$, $b=\code h$ (so $\acl(a)=\acl(g)$ etc.).  Since $\Cb(hg^{-1})\subseteq\acl(g,h)$,
$\Cb(h)\subseteq\acl(c,g)$ ($h=cg$) and $\Cb(g)\subseteq\acl(c,h)$ ($g=c^{-1}h$), \ref{GC:group}
holds.  $x=g(y)\in\dcl(g,y)$ and $y=g^{-1}(x)$ ($x\in\operatorname{Dom}(g^{-1})$ because
$x$ is generic over $\acl(\Cb g)$, \cref{lem:eval}(2)), so $\acl(ax)=\acl(ay)$; similarly
$\acl(by)=\acl(bz)$.  Next, $x\ind{}gh$: $y\ind{}gh$ and $x\interalg y$ over $g$ give
$\rk(x/gh)=\rk(y/gh)=m$.  Hence $x\in\operatorname{Dom}(c)$ and by \cref{lem:comp-codes}(2)
$c(x)=h(g^{-1}(x))=h(y)=z$; so $z\in\dcl(c,x)$, $x=c^{-1}(z)$ and $\acl(cx)=\acl(cz)$.
This proves \ref{GC:space}.

\ref{GC:cb}: by \cref{lem:triple} the germ determined by the triple $(a;y,x)$ (with
$y\ind{}a$ by hypothesis, $x\ind{}a$ by \cref{lem:genset-basic}(2), $x\in\dcl(a,y)$ by
\cref{lem:eval-dcl}) is the unique germ with code in $\acl(a)$ sending $y$ to $x$, hence
is $g$ itself; so $\Cb(yx/a)=\Cb(g)=k(a)$ and, as $\FOD$ is invariant under renaming the
variables (\cref{thm:weil-fod}), $\Cb(xy/a)=\Cb(yx/a)$, whence
$a\subseteq\acl(\Cb(xy/a))$.  Likewise for $b$ with $(y,z)$ and for $c$ with $(x,z)$
(using $x\ind{}c$, $z\ind{}c$, which follow from $\rk(x/gh)=m$ and
$\rk(z/gh)=\rk(y/gh)=m$).

\ref{GC:indep}: all six points lie in $\acl(ghy)$, which has rank $2n+m$.  We check the
$16$ non-collinear triples.  Since $\acl(ab)=\acl(ac)=\acl(bc)$ it suffices, for triples
with two group points, to treat $(g,h,p)$ with $p\in\{x,y,z\}$: $\rk(ghy)=2n+m$ by
hypothesis, and $\rk(x/gh)=\rk(z/gh)=m$ was shown above.  Triples with one group point
and two space points not on its line: $(g,y,z)$, $(g,x,z)$, $(h,x,y)$, $(h,x,z)$,
$(c,x,y)$, $(c,y,z)$.  From $g\ind{}hy$ and $z\in\acl(hy)$ we get $g\ind{}yz$, so
$\rk(z/gy)=\rk(yz/g)-\rk(y/g)=\rk(yz)-m=m$, using $\rk(yz)=2m$ (which is \ref{H:trans}
for $h$); hence $\rk(gyz)=n+2m$, and $\rk(gxz)=\rk(gyz)$ as $\acl(gx)=\acl(gy)$.
Symmetrically $h\ind{}gy$ gives $\rk(hxy)=\rk(hxz)=n+2m$.  For $c$: $c$ and $h$ are
interalgebraic over $g$, so $\rk(c/xyg)=\rk(h/xyg)=\rk(h/yg)=n$ (as $x\in\acl(gy)$);
therefore $\rk(c/xy)=n$, $c\ind{}xy$ and $\rk(cxy)=n+2m$; similarly $c$ and $g$ are
interalgebraic over $h$, $\rk(c/yzh)=\rk(g/yzh)=\rk(g/yh)=n$, so $\rk(cyz)=n+2m$.
Finally $(x,y,z)$: $\rk(xyz)=\rk(y)+\rk(x/y)+\rk(z/xy)$ with $\rk(x/y)=m$ (\ref{H:trans})
and $\rk(z/xy)=\rk(z/xyg)+\bigl(\rk(g/xy)-\rk(g/xyz)\bigr)$; here $\rk(z/xyg)=\rk(z/yg)=m$
(shown above: $g\ind{}hy$ and $z\in\acl(hy)$ give $g\ind{}yz$, so $\rk(z/gy)=m$) and $\rk(g/xy)\ge\rk(g/xyz)$, so
$\rk(z/xy)\ge m$, hence $=m$ and $\rk(xyz)=3m$.  \ref{GC:rank} is $n\ge1$, which holds
since $\mathcal G$ consists of germs of rank $n$ moving generic points
(\ref{H:trans} with $m\ge1$ forces $n\ge1$: a germ of rank $0$ is $\Gamma$-invariant and
$g(y)\in\dcl(k,y)$ would then be interalgebraic with $y$, contradicting $g(y)\ind{}y$).
\end{proof}
